\documentclass[11pt]{article}

%-------------------------
% Packages
%-------------------------
\usepackage[a4paper,margin=1in]{geometry}
\usepackage{xcolor}
\usepackage{tikz}
\usepackage[most]{tcolorbox}
\usepackage{mathtools}
\usepackage{amsthm, amsmath, amssymb}
\usepackage{enumitem}
\usepackage{hyperref}
\usepackage[nameinlink,noabbrev]{cleveref}
\usepackage{titling} 
\usepackage{fontspec}
\usepackage{unicode-math}


%-------------------------
% Fonts
%-------------------------
\setmainfont{JetBrains Mono}
\setsansfont{JetBrains Mono}
\setmonofont{JetBrains Mono}


\setmathfont[
    Scale=1.25
]{latinmodern-math.otf}

\setmathfont[
  range=\mathup/{latin, Latin, 0-9},
  range=\mathit/{latin, Latin},
  range=\mathbfit/{latin, Latin},
  Scale=1.5
]{JetBrains Mono Bold} 

\setmathfont[range=\mathbb, Scale=MatchUppercase]{latinmodern-math.otf}

%-------------------------
% Paragraph spacing
%-------------------------
\setlength{\parindent}{0pt}
\setlength{\parskip}{0.7em}


%-------------------------
% Dark mode colors
%-------------------------
\definecolor{bgcolor}{HTML}{333333}
\definecolor{textcolor}{HTML}{FAFAFA}
\definecolor{defcolor}{HTML}{E86873}
\definecolor{thmcolor}{HTML}{0A9396}
\definecolor{lemcolor}{HTML}{94D2BD}
\definecolor{corcolor}{HTML}{9B4AF7}
\definecolor{probcolor}{HTML}{EE9B00}
\definecolor{excolor}{HTML}{21E933}
\definecolor{propcolor}{HTML}{EE9B20}

%-------------------------
% Background and text color
%-------------------------
\pagecolor{bgcolor}
\color{textcolor}

%-------------------------
% Paragraph spacing
%-------------------------
\setlength{\parindent}{0pt}
\setlength{\parskip}{0.7em}

%-------------------------
% Base tcolorbox settings (Removed 'attach title to upper')
%-------------------------
\tcbset{
  enhanced,
  colback=bgcolor,
  colframe=thmcolor,
  coltext=white,
  coltitle=white,
  fonttitle=\bfseries,
  boxrule=0.7pt,
  left=1em,
  right=1em,
  top=0.7em,
  bottom=0.7em,
  before skip=10pt,
  after skip=10pt,
  title after break,
  breakable
}

%-------------------------
% Standard Theorem Environments (using amsthm)
% Syntax: \begin{thm}[Optional Title]\label{thm:label} ...
%-------------------------
\newtheorem{thm}{Theorem}[section]
\newtheorem*{thm*}{Theorem}
\newtheorem{defn}{Definition}[section]
\newtheorem*{defn*}{Definition}
\newtheorem{lem}{Lemma}[section]
\newtheorem*{lem*}{Lemma}
\newtheorem{cor}{Corollary}[section]
\newtheorem*{cor*}{Corollary}
\newtheorem{prob}{Problem}[section]
\newtheorem*{prob*}{Problem}
\newtheorem{ex}{Example}[section]
\newtheorem*{ex*}{Example}
\newtheorem{prop}{Proposition}[section]
\newtheorem*{prop*}{Proposition}

%-------------------------
% Tcolorbox environments (applies style to amsthm environments)
% Removed 'attach title to upper' from all individual settings.
%-------------------------
\tcolorboxenvironment{thm}{
  colframe=thmcolor,
  colback=thmcolor!15!bgcolor,
  fonttitle=\bfseries,
  coltitle=white,
  breakable,
  title after break,
}

\tcolorboxenvironment{thm*}{
  colframe=thmcolor,
  colback=thmcolor!15!bgcolor,
  fonttitle=\bfseries,
  coltitle=white,
  breakable,
  title after break,
}

\tcolorboxenvironment{defn}{
  colframe=defcolor,
  colback=defcolor!15!bgcolor,
  fonttitle=\bfseries,
  coltitle=white,
  breakable,
  title after break,
}

\tcolorboxenvironment{defn*}{
  colframe=defcolor,
  colback=defcolor!15!bgcolor,
  fonttitle=\bfseries,
  coltitle=white,
  breakable,
  title after break,
}

\tcolorboxenvironment{lem}{
  colframe=lemcolor,
  colback=lemcolor!15!bgcolor,
  fonttitle=\bfseries,
  coltitle=white,
  breakable,
  title after break,
}

\tcolorboxenvironment{lem*}{
  colframe=lemcolor,
  colback=lemcolor!15!bgcolor,
  fonttitle=\bfseries,
  coltitle=white,
  breakable,
  title after break,
}

\tcolorboxenvironment{cor}{
  colframe=corcolor,
  colback=corcolor!15!bgcolor,
  fonttitle=\bfseries,
  coltitle=white,
  breakable,
  title after break,
}

\tcolorboxenvironment{cor*}{
  colframe=corcolor,
  colback=corcolor!15!bgcolor,
  fonttitle=\bfseries,
  coltitle=white,
  breakable,
  title after break,
}

\tcolorboxenvironment{prob}{
  colframe=probcolor,
  colback=probcolor!15!bgcolor,
  fonttitle=\bfseries,
  coltitle=white,
  breakable,
  title after break,
}

\tcolorboxenvironment{prob*}{
  colframe=probcolor,
  colback=probcolor!15!bgcolor,
  fonttitle=\bfseries,
  coltitle=white,
  breakable,
  title after break,
}

\tcolorboxenvironment{ex}{
  colframe=excolor,
  colback=excolor!15!bgcolor,
  fonttitle=\bfseries,
  coltitle=white,
  breakable,
  title after break,
}

\tcolorboxenvironment{ex*}{
  colframe=excolor,
  colback=excolor!15!bgcolor,
  fonttitle=\bfseries,
  coltitle=white,
  breakable,
  title after break,
}

\tcolorboxenvironment{prop}{
    colframe=propcolor, 
    colback=propcolor!15!bgcolor,
    fonttitle=\bfseries,
    coltitle=white,
    breakable,
    title after break
}

\tcolorboxenvironment{prop*}{
    colframe=propcolor, 
    colback=propcolor!15!bgcolor,
    fonttitle=\bfseries,
    coltitle=white,
    breakable,
    title after break
}

%-------------------------
% Cleveref settings (Using short environment names for correct display)
%-------------------------
\crefname{thm}{theorem}{theorems}
\Crefname{thm}{Theorem}{Theorems}

\crefname{defn}{definition}{definitions}
\Crefname{defn}{Definition}{Definitions}

\crefname{lem}{lemma}{lemmas}
\Crefname{lem}{Lemma}{Lemmas}

\crefname{cor}{corollary}{corollaries}
\Crefname{cor}{Corollary}{Corollaries}

\crefname{prob}{problem}{problems}
\Crefname{prob}{Problem}{Problems}

\crefname{ex}{example}{examples}
\Crefname{ex}{Example}{Examples}

\crefname{prop}{proposition}{propositions}
\Crefname{prop}{Proposition}{Propositions}

\crefname{thm*}{theorem}{theorems}
\Crefname{thm*}{Theorem}{Theorems}

\crefname{defn*}{definition}{definitions}
\Crefname{defn*}{Definition}{Definitions}

\crefname{lem*}{lemma}{lemmas}
\Crefname{lem*}{Lemma}{Lemmas}

\crefname{cor*}{corollary}{corollaries}
\Crefname{cor*}{Corollary}{Corollaries}

\crefname{prob*}{problem}{problems}
\Crefname{prob*}{Problem}{Problems}

\crefname{ex*}{example}{examples}
\Crefname{ex*}{Example}{Examples}

\crefname{prop*}{proposition}{propositions}
\Crefname{prop*}{Proposition}{Propositions}


\title{\huge{Assignment 3}}
\author{\LARGE{Remote grp}}
\date{for September 29, 2026}

\begin{document}
\maketitle


\begin{prob*}[Munkres 16.9]
    Show that 
    the dictionary order topology on the set 
    $\mathbb R \times \mathbb R$ is the 
    same as the product topology $\mathbb R_d 
    \times \mathbb R$ where $\mathbb R_d$ renotes 
    $\mathbb R$ in the discrete topology. 
    Compare this topology with the standard 
    topology on $\mathbb R^2$. 
\end{prob*}

\begin{proof}
    Recall that the dictionary order on $\mathbb R^2$ is defined by
    \[
        (x_1,x_2)\prec (y_1,y_2)
    \]
    if either $x_1<y_1$, or $x_1=y_1$ and $x_2<y_2$.

    Let $\mathcal T_{\mathrm{dict}}$ denote the dictionary order topology
    and let $\mathcal T_{\mathrm{prod}}$ denote the product topology on
    $\mathbb R_d\times \mathbb R$.

    Since $\mathbb R_d$ is discrete, a basis for
    $\mathcal T_{\mathrm{prod}}$ is given by sets of the form
    \[
        \{x\}\times (a,b),
    \]
    but
    \[
        \{x\}\times(a,b)
        =
        ((x,a),(x,b))
    \]
    in the dictionary order. Hence every basis element of
    $\mathcal T_{\mathrm{prod}}$ is open in
    $\mathcal T_{\mathrm{dict}}$. Therefore
    \[
        \mathcal T_{\mathrm{prod}}
        \subseteq
        \mathcal T_{\mathrm{dict}}.
    \]

    Conversely, consider a basic open interval
    \[
        ((a,b),(c,d))
    \]
    in the dictionary order.

    If $a=c$, then
    \[
        ((a,b),(a,d))
        =
        \{a\}\times(b,d),
    \]
    which is open in $\mathbb R_d\times\mathbb R$.

    If $a<c$, then
    \[
        ((a,b),(c,d))
        =
        \bigl(\{a\}\times(b,\infty)\bigr)
        \cup
        \bigl((a,c)\times\mathbb R\bigr)
        \cup
        \bigl(\{c\}\times(-\infty,d)\bigr).
    \]
    Since $\mathbb R_d$ is discrete, every subset of $\mathbb R_d$
    is open, so each of these sets is open in
    $\mathbb R_d\times\mathbb R$.

    Thus every basic open set in the dictionary order topology is open
    in the product topology, and hence
    \[
        \mathcal T_{\mathrm{dict}}
        \subseteq
        \mathcal T_{\mathrm{prod}}.
    \]

    Therefore
    \[
        \mathcal T_{\mathrm{dict}}
        =
        \mathcal T_{\mathrm{prod}}.
    \]

    Finally, we compare this with the usual topology on $\mathbb R^2$.
    Since the discrete topology on $\mathbb R$ is finer than the usual
    topology,
    \[
        \mathbb R_d\times\mathbb R
    \]
    has a topology finer than the usual product topology on $\mathbb R^2$.

    The inclusion is strict. For example,
    \[
        \{0\}\times(-1,1)
    \]
    is open in $\mathbb R_d\times\mathbb R$, hence in the dictionary
    order topology, but it is not open in the usual topology on
    $\mathbb R^2$.

    Hence the dictionary order topology is strictly finer than the
    standard topology on $\mathbb R^2$.
\end{proof}


\end{document}
